\chapter{1944:Otto Hahn}

\label{chap:chemistry1944}

The Nobel Prize in Chemistry for the year 1944 was awarded to Otto Hahn.

\textbf{Motivation:}

for his discovery of the fission of heavy nuclei

\section{Otto Hahn}

\subsection*{Early Life and Education}

Otto Hahn was born on 1879-03-08 in Frankfurt am Main, Germany.
He was a German radiochemist awarded the Nobel Prize in Chemistry in 1944 for his discovery of the fission of heavy nuclei.

\subsection*{Career and Major Research}

Hahn studied chemistry at the University of Marburg, where he received his doctorate in 1901 under Theodor Zincke.
He worked with William Ramsay in London, where he discovered radiothorium in 1905, and with Ernest Rutherford in Montreal from 1905 to 1906.
From 1912 he headed the radioactivity department of the Kaiser Wilhelm Institute for Chemistry in Berlin-Dahlem, where he worked for many years with Lise Meitner and later with Fritz Strassmann.
He discovered mesothorium and, together with Meitner, the element protactinium in 1917.
In 1938, together with Fritz Strassmann, he discovered nuclear fission: bombarding uranium with neutrons, they detected barium among the reaction products and showed that the uranium nucleus had split into lighter nuclei.

\subsection*{Nobel Prize-winning Work}

The work for which Otto Hahn was awarded the Nobel Prize in Chemistry in 1944 was described by the Nobel Committee as: "for his discovery of the fission of heavy nuclei".

The chemical evidence that uranium, on bombardment with neutrons, is converted into lighter elements such as barium established that the atomic nucleus can be split.
The process was immediately explained theoretically by Lise Meitner and Otto Frisch, who gave it the name fission.

\subsection*{Significance and Impact}

The discovery of nuclear fission changed the course of history.
It made possible both the atomic bomb and the peaceful use of nuclear energy, and it earned Hahn the 1944 Nobel Prize in Chemistry, which was presented to him after the end of the Second World War.

\subsection*{Later Career and Honors}

After the war Hahn was president of the Kaiser Wilhelm Society and, from 1946, of its successor the Max Planck Society, serving until 1960.
He received the Enrico Fermi Award together with Meitner and Strassmann in 1966, and he was an outspoken opponent of nuclear weapons.
He died on 1968-07-28 in Göttingen, Germany.

\subsection*{Legacy}

The legacy of Otto Hahn endures through the lasting impact of his scientific contributions.
The discovery of nuclear fission, made in his laboratory, remains one of the defining events of twentieth-century science.

\subsection*{Selected Publications}

\begin{itemize}

\item Hahn, O., \& Strassmann, F. (1939). "Über den Nachweis und das Verhalten der bei der Bestrahlung des Urans mittels Neutronen entstehenden Erdalkalimetalle." Naturwissenschaften, 27(1), 11--15.

\item Hahn, O. (1962). "Vom Radiothor zur Uranspaltung." Braunschweig: Vieweg.

\end{itemize}

\clearpage

\section*{Chapter Summary}

The 1944 prize recognized Otto Hahn for his discovery of the fission of heavy nuclei. The 1938 discovery of nuclear fission, with Fritz Strassmann, opened the nuclear age and led both to nuclear weapons and to nuclear energy.
